\chapter{feasibility}
\section{Technical Feasibility (High Viability)}
The technical risk profile for the MotionBridge framework is exceptionally low due to the structural selection of mature, optimized, and heavily documented open-source ecosystems.

\begin{itemize}
    \item \textbf{Software Ecosystem Synergy:} The pipeline leverages standard Google \texttt{MediaPipe} libraries (Pose and Hands Landmarkers) for lightweight 2D frame feature extraction, \texttt{OpenCV} for spatial coordinate triangulation math, and \texttt{PyBullet} for forward/inverse kinematic simulation mapping. These packages possess proven, stable interoperability in standard Python 3.9+ execution environments.
    \item \textbf{Hardware Democratic Accessibility:} Traditional optical or inertial motion capture topologies rely on high-overhead dedicated GPU architecture or high-cost active depth configurations. MotionBridge bypasses these barriers entirely; both \texttt{MediaPipe} and \texttt{PyBullet} are natively optimized for concurrent Central Processing Unit (CPU) thread loops. Standard Intel Core i5 (8th Generation) host machines with 8 GB RAM satisfy performance thresholds.
    \item \textbf{Algorithmic Resilience to Noise:} The system natively counteracts regional body self-occlusion through a custom confidence-weighted triangulation algorithm:
    \begin{equation}
        X_{\text{3D}} = \frac{\sum_{i=1}^{3} c_i \cdot f(P_i, x_i)}{\sum_{i=1}^{3} c_i}
    \end{equation}
    By dynamically discounting occluded viewpoints where the landmark visibility score ($c_i \to 0$), geometric baseline errors are mitigated. The proposal to introduce analytical temporal smoothing (\textit{One Euro Filter}) dynamically scaled via:
    \begin{equation}
        \alpha = \frac{1}{1 + \tau \Delta t}, \quad \tau = \frac{1}{2\pi \cdot (f_{\text{min}} + \beta \cdot |\dot{X}_{\text{3D}}|)}
    \end{equation}
    further isolates high-frequency simulation jitter from propagating into the downstream trajectory tables.
\end{itemize}

\section{Economic Feasibility (Excellent Viability)}
From a budgetary and resource-allocation perspective, the MotionBridge system exhibits exceptional cost efficiency, rendering it highly feasible for deployment within low-income research environments and academic laboratories.

\begin{itemize}
    \item \textbf{Capital Asset Minimization:} Traditional industry standard spatial marker systems cost thousands of dollars. The entire physical equipment specification for MotionBridge scales under an estimated total budget of \textbf{NPR 24,000} (approximately \$180 USD), as detailed below:
\end{itemize}

\begin{table}[h!]
\centering
\caption{Estimated Hardware and Project Budget Breakdown}
\begin{tabular}{llrr}
\toprule
\textbf{S.N.} & \textbf{Item Description} & \textbf{Qty} & \textbf{Total Price (NPR)} \\
\midrule
1 & Standard 1080p USB Webcams & 3 & 12,000 \\
2 & Fixed Camera Tripods & 3 & 4,500 \\
3 & Powered USB Extension Hub & 1 & 2,000 \\
4 & Geometric Checkerboard Calibration Target & 1 & 1,000 \\
5 & Open-Source Core Software Modules & - & Free \\
6 & Technical Project Documentation \& Testing & - & 2,500 \\
7 & Operational Contingency Reserve Fund & - & 2,000 \\
\midrule
\textbf{} & \textbf{Total Estimated Budget} & & \textbf{NPR 24,000} \\
\bottomrule
\end{tabular}
\end{table}

\begin{itemize}
    \item \textbf{Zero Software Overhead:} The complete computational development stack (\texttt{Python}, \texttt{OpenCV}, \texttt{MediaPipe}, \texttt{PyBullet}, \texttt{Pandas}, \texttt{Git}) contains zero proprietary licensing or execution recurring fees.
\end{itemize}

\section{Operational Feasibility (High Viability)}
Operational feasibility evaluates how effectively the system integrates into actual deployment settings to produce the desired outcomes.

\begin{itemize}
    \item \textbf{Strategic Scope Bounding:} The team has managed operational risk by structurally separating data collection from simulation constraints. While intricate 21-point hand finger joint coordinates are logged to the output files, they are explicitly omitted from the \texttt{PyBullet} simulation visualization due to standard Kuka IIWA URDF link structural limitations. This boundary preserves execution frame rates.
    \item \textbf{High-Value Extensible Artifact Production:} The pipeline runs an automated data factory serialization loop that seamlessly exports structured CSV parameters:
    \begin{equation}
        \text{Record} = \{t, X_{\text{shoulder}}, X_{\text{elbow}}, X_{\text{wrist}}, X_{\text{fingers}}, \theta\}
    \end{equation}
    This highly clean, structured timeline format is directly compatible with downstream imitation learning pipelines, behavioral cloning neural models, or inverse reinforcement learning protocols.
    \item \textbf{Environmental Adaptability:} The architecture operates in any standard ambient space. It requires only an initial checkerboard coordinate matrix execution pass ($9 \times 6$ corners) and consistent indoor lighting to verify structural integrity.
\end{itemize}

\section{Schedule Feasibility (Realistic Viability)}
The project layout presents a sustainable and logically organized developmental progression over the scheduled timeline:
\begin{itemize}
    \item Structural phases are explicitly partitioned into upfront research and literature consolidation, reducing design gaps.
    \item A massive mid-timeline block is dedicated purely to core \textbf{System Design and Development}, allowing the team sufficient cushion to iron out multi-threaded asynchronous frame sync loops ($\Delta t < 15\text{ ms}$).
    \item Adequate downstream space is reserved for system validation, constraint testing under simulated occlusions, and documentation refinement prior to final submittal.
\end{itemize}